%-----------------------------------------------------------------------------
%  Aerospace-Module-2 - Relazione di progetto
%
%  Compilazione:  make        (oppure: latexmk -pdf main.tex)
%  Le figure dei risultati non sono duplicate qui dentro: \graphicspath punta
%  direttamente alle cartelle del progetto che le generano.
%-----------------------------------------------------------------------------
\documentclass[11pt,a4paper]{article}

% --- lingua e font ---------------------------------------------------------
\usepackage[T1]{fontenc}
\usepackage[utf8]{inputenc}
\usepackage{lmodern}
\usepackage[italian]{babel}
\usepackage{microtype}
\usepackage{csquotes}

% --- pagina ----------------------------------------------------------------
\usepackage[margin=2.7cm]{geometry}
\usepackage{fancyhdr}

% --- matematica, tabelle, figure -------------------------------------------
\usepackage{amsmath}
\usepackage{amssymb}
\usepackage{graphicx}
\usepackage{booktabs}
\usepackage{float}
\usepackage{subcaption}
\usepackage{xcolor}
\usepackage{enumitem}

% --- codice ----------------------------------------------------------------
\usepackage{listings}
\usepackage{upquote}

% --- bibliografia ----------------------------------------------------------
\usepackage[numbers,sort&compress]{natbib}

% --- collegamenti (per ultimo) ---------------------------------------------
\usepackage[hidelinks]{hyperref}

%-----------------------------------------------------------------------------
%  Impostazioni
%-----------------------------------------------------------------------------

% figure prese dove il progetto le produce, senza copiarle nel report
\graphicspath{%
  {immagini/}%
  {../applications/test/nonEqTTv/output/}%
  {../explainations/dipendenze/}%
}

\pagestyle{fancy}
\fancyhf{}
\fancyhead[L]{\small\itshape Modello a due temperature in OpenFOAM-13 con Mutation++}
\fancyhead[R]{\small\thepage}
\renewcommand{\headrulewidth}{0.4pt}

\definecolor{codebg}{gray}{0.96}
\definecolor{codekw}{rgb}{0.15,0.30,0.60}
\definecolor{codecm}{rgb}{0.35,0.50,0.35}

\lstdefinestyle{progetto}{%
  backgroundcolor=\color{codebg},
  basicstyle=\ttfamily\footnotesize,
  keywordstyle=\color{codekw}\bfseries,
  commentstyle=\color{codecm}\itshape,
  breaklines=true,
  showstringspaces=false,
  columns=fullflexible,
  frame=single,
  framerule=0pt,
  xleftmargin=0.6em,
  framexleftmargin=0.6em,
  captionpos=b,
}
\lstset{style=progetto}

% notazione ricorrente
\newcommand{\Ttr}{T_{tr}}
\newcommand{\Tve}{T_{ve}}
\newcommand{\eve}{e_{ve}}
\newcommand{\hyfoam}{\texttt{hy2Foam}}
\newcommand{\mpp}{Mutation\texttt{++}}
\newcommand{\shockThermo}{\texttt{shockThermo}}
\newcommand{\file}[1]{\texttt{#1}}

% segnaposto per le sezioni ancora da scrivere
\newcommand{\daScrivere}{%
  \par\medskip\noindent
  \textcolor{gray}{\itshape [sezione da scrivere]}\par\medskip
}

%-----------------------------------------------------------------------------
\begin{document}

\begin{titlepage}
\centering
\vspace*{2cm}

{\Large Politecnico di Milano\par}
{\large Dipartimento di Scienze e Tecnologie Aerospaziali\par}

\vfill

{\LARGE\bfseries Un modello a due temperature\\[0.2em]
per flussi ipersonici reagenti\\[0.4em]
in OpenFOAM-13 e \mpp\par}

\vspace{1.2cm}

{\large Riproduzione degli heat bath zero-dimensionali di\\
Casseau \textit{et al.}, \textit{Aerospace} \textbf{3} (2016) 34\par}

\vfill

{\large Daniele Salvi\par}

\vspace{0.8cm}

{\large \today\par}

\vspace*{2cm}
\end{titlepage}

\section*{Sommario}
\addcontentsline{toc}{section}{Sommario}

Questa relazione documenta la riproduzione, in OpenFOAM-13 e con la libreria
termochimica \mpp, dei casi di verifica zero-dimensionali con cui Casseau
\textit{et al.}~\cite{casseau2016} hanno validato il solver a due temperature
\hyfoam. Il modello di Park a due temperature -- una temperatura
traslazionale-rotazionale $\Ttr$ e una vibro-elettronica $\Tve$ -- è stato
implementato in forma conservativa, risolvendo le equazioni nelle energie
anziché nelle temperature e delegando a \mpp{} sia l'inversione delle energie
in temperature sia il calcolo delle proprietà termodinamiche.

Il lavoro procede su due percorsi che condividono le stesse funzioni sorgente:
un insieme di programmi zero-dimensionali autonomi, che integrano le equazioni
del bagno termico chiamando soltanto \mpp{} e coprono tutte le figure del
capitolo di verifica del paper, e il solver CFD, che risolve le stesse equazioni su una singola cella
attraverso un modello termodinamico che fa da ponte verso la libreria. L'accordo fra i due percorsi, entro lo 0{,}03--0{,}2\%
su tutte le curve, verifica l'integrazione del modello nel solver; il confronto
con le curve di \hyfoam{} estratte dalle figure vettoriali del paper e con le
temperature di equilibrio ricavate dalla conservazione dell'energia ne verifica
la fisica.


\tableofcontents
\clearpage

\section{Introduzione}
\label{sec:introduzione}

\subsection{Il non-equilibrio termochimico nel rientro atmosferico}
\label{sec:contesto}

Un veicolo che rientra nell'atmosfera a velocità ipersonica genera un urto
forte, dietro il quale la temperatura sale a decine di migliaia di kelvin: i
modi vibrazionali si eccitano, i livelli elettronici si popolano e le molecole
si dissociano. Questi processi non raggiungono l'equilibrio alla stessa
velocità: traslazione e rotazione si equilibrano in pochi urti molecolari, la
vibrazione in molti di più, e il gas dietro l'urto si trova in non-equilibrio
termico. Il modello a due temperature di Park~\cite{park1990,park1993} lo
descrive tenendo distinte una temperatura traslazionale-rotazionale e una
vibro-elettronica, legate da un tempo di rilassamento. A questo si somma
l'accoppiamento con la chimica: una molecola eccitata vibrazionalmente si
dissocia più facilmente, e le reazioni a loro volta sottraggono o restituiscono
energia vibrazionale. Dalla descrizione di questi processi dipende la
previsione dei carichi termici sul veicolo.

\subsection{Il modello a una temperatura e i suoi limiti}
\label{sec:una-temperatura}

Una molecola può immagazzinare energia in modi diversi: spostandosi nello
spazio (traslazione), ruotando su se stessa (rotazione), facendo oscillare i
propri atomi (vibrazione) ed eccitando i propri elettroni. Nella maggior parte
delle simulazioni CFD, compresi i solver standard di OpenFOAM, il gas è
descritto da una sola temperatura: si assume che l'energia ricevuta dal gas si
ripartisca subito fra tutti questi modi nelle proporzioni di equilibrio, per
cui un solo numero basta a sapere quanta energia c'è in ciascuno, e con lo
stesso numero si calcola anche la velocità delle reazioni chimiche. È
un'ipotesi eccellente nella gran parte delle applicazioni, perché le molecole
si urtano così spesso che la ripartizione avviene molto più in fretta di
quanto cambi il flusso.

Dietro un urto ipersonico questa ipotesi cade: l'energia arriva prima al moto
delle molecole e solo dopo molti urti alla vibrazione, in un tempo paragonabile
a quello che il gas impiega ad attraversare la regione dietro l'urto. Un
modello a una temperatura non vede questo ritardo: prevede appena dietro l'urto
una temperatura più bassa di quella reale e salta la zona in cui il gas si
porta gradualmente all'equilibrio. Lo stesso accade per la chimica: la
dissociazione dipende sia dall'energia con cui le molecole si urtano sia da
quanto stanno già vibrando, due grandezze che una sola temperatura non può
distinguere. Il risultato sono errori sulla temperatura e sulla composizione
del gas, e quindi sulla posizione dell'urto e sul calore che raggiunge la
parete~\cite{park1990}.

\subsection{La soluzione: il modello a due temperature}
\label{sec:soluzione}

Il modello a due temperature di Park~\cite{park1990,park1993} supera questi
limiti usando due temperature invece di una: la prima descrive il moto e la
rotazione delle molecole, che si equilibrano quasi subito, la seconda la
vibrazione e gli elettroni, che restano indietro. Subito dietro l'urto le due
possono differire di migliaia di gradi; l'energia passa poi dall'una all'altra
gradualmente, tanto più in fretta quanto più il gas è caldo e denso, finché le
due temperature coincidono e il gas è tornato in equilibrio. Anche la chimica
le usa entrambe: la dissociazione è calcolata a una temperatura intermedia,
così che una molecola che vibra ancora poco si dissoci più lentamente, e
l'energia di vibrazione portata via dalle molecole che si spezzano viene tolta
da quella del gas. Resta un'approssimazione, perché dentro ciascuno dei due
gruppi l'energia è supposta in equilibrio, ma cattura l'essenziale senza dover
seguire ogni singolo livello energetico delle molecole.

\subsection{Il lavoro di riferimento}
\label{sec:riferimento}

Il lavoro di riferimento è l'articolo di Casseau
\textit{et al.}~\cite{casseau2016}, dell'Università di Strathclyde, che
presenta \hyfoam, un solver a due temperature per flussi ipersonici reagenti
costruito su OpenFOAM e distribuito come codice aperto~\cite{hystrath}. La
prima parte dell'articolo verifica il solver su una serie di bagni termici:
piccoli volumi di gas fermi e isolati, preparati fuori equilibrio e lasciati
evolvere, in cui ogni processo si può studiare da solo. I casi vanno dall'azoto
puro, con e senza energia elettronica, alle miscele di azoto con azoto atomico
e con ossigeno, fino ai gas che reagiscono e all'aria a cinque specie, e i
risultati sono confrontati con altri codici CFD e con codici DSMC, che
simulano il gas seguendo il moto e gli urti di un gran numero di particelle.
Sono questi i casi che il presente lavoro riproduce.

\clearpage

\section{Il modello a due temperature}
\label{sec:modello}

\daScrivere

% Traccia:
%  - decomposizione dell'energia nei modi (eq. 1-7 del paper) e le due T
%  - variabili conservate: rho_s, E, E_ve (eq. 23); perche' energie e non T
%  - equazione di stato: p = somma delle pressioni parziali (eq. 24)
%  - scambio V-T: Landau-Teller (eq. 8), tau di Millikan-White + Park (eq. 9-17)
%  - scambio V-V: Knab (eq. 18), solo nel caso N2-O2
%  - chimica: legge di azione di massa (eq. 27), Arrhenius (eq. 28),
%    temperatura di Park T^a Tv^(1-a) (eq. 29)
%  - accoppiamento chimica-vibrazione: modello non preferenziale (eq. 30-31)
%  - riepilogo del sistema risolto nel bagno termico (eq. 22, 25, 26)

\clearpage

\section{Implementazione}
\label{sec:implementazione}

\daScrivere

% Traccia:
%  - quadro d'insieme: i due percorsi e l'header condiviso mutationSources.H
%  - il thermo come ponte: gerarchia OpenFOAM-13
%    (HighEnthalpyMulticomponentThermo -> MulticomponentThermo -> PsiThermo
%     -> BasicThermo<MixtureType, composite>), i campi Tve_ ed eve_
%  - il ciclo cella per cella: setState con densita' parziali e temperature
%    (vars=1), con energie (vars=0) per la decodifica, con p,T,Tve (vars=2)
%    per psi
%  - le tre sorgenti esposte al solver: computeSourceVe, computeSourceY,
%    computeSourceE
%  - il solver: thermophysicalPredictor (YiEqn, EEqn, EveEqn, correct) e
%    perche' l'EEqn e' ricopiata invece di chiamare shockFluid
%  - shockFluid: copia del solo header del modulo core, fluidThermo al posto
%    di psiThermo e costruttore che riceve il thermo gia' costruito
%  - i programmi 0D e heatBath.H (energie conservate, T di equilibrio)
%  - i dati: cosa e' del progetto (species.xml, meccanismi) e cosa e' default
%  - catena di compilazione ed esecuzione (Allwmake, Allrun, MPP_DATA_DIRECTORY)
%  - figura: pannelli delle dipendenze in explainations/dipendenze/

\clearpage

\section{Casi di verifica e risultati}
\label{sec:validazione}

\daScrivere

% Traccia:
%  - metodo del confronto a tre vie: solver vs 0D, 0D vs paper, T di equilibrio
%  - estrazione delle curve dal pdf vettoriale e sua incertezza (40-60 K)
%  - tabella riassuntiva solver vs 0D (0.03-0.2 %)
%  - tabella riassuntiva 0D vs paper (T finali e scarti nel transitorio)
%  - figura per figura: 3a, 3b, 4 (senza/con E_el), 5, 6 (senza/con V-V),
%    7, 8, 9 (Park e QK)
%  - grafici da output/figX.png e output/figX-n.png
%  - costo di calcolo: solver ~28 min contro ~2 s del programma 0D

\clearpage

\section{Scelte di modellazione}
\label{sec:discussione}

\daScrivere

% Traccia:
%  - forza motrice del V-T: e_ve = e_v + e_el come nel paper, non l'OmegaVT
%    di Mutation++ (28 %/56 % -> 0.3 %/4 % sulla fig. 4 con E_el)
%  - energia elettronica accesa solo per la fig. 4: le T di equilibrio delle
%    fig. 5-8 tornano solo senza; il testo del paper dice il contrario della
%    figura
%  - media dei tempi di rilassamento: aritmetica in Mutation++, armonica per
%    coppie nel paper; e' l'origine del ritardo su Tv nelle miscele
%  - esponente della temperatura di Park: 0.7 (paper) contro 0.5 (Mutation++),
%    effetto sulla composizione della fig. 7
%  - accoppiamento chimica-vibrazione non preferenziale
%  - fig. 9: costanti QK invece di Park, e perche'
%  - sezione d'urto V-V di N2-O2 presa da hyStrath

\clearpage

\section{Limiti e sviluppi}
\label{sec:limiti}

\daScrivere

% Traccia:
%  - l'equazione di eve e' risolta senza trasporto: va bene per il bagno
%    termico, non per un caso con flusso
%  - proprieta' di trasporto ancora dal thermo base; casi validati inviscidi
%  - il thermo base e' valido fino a 20000 K e lo segnala
%  - costo: le velocita' di reazione sono ricalcolate piu' volte per cella
%  - densita' usata nelle sorgenti (psi*p) sfasata di un sotto-passo PIMPLE
%  - condizioni al contorno del campo eve
%  - sviluppi: trasporto e diffusione di eve, tubo d'urto 1D, casi 2D/3D

\clearpage

\section{Conclusioni}
\label{sec:conclusioni}

\daScrivere

\clearpage

\bibliographystyle{unsrtnat}
\bibliography{bibliografia}

\end{document}
